\documentclass[11pt]{article}
\usepackage[a4paper,margin=1in]{geometry}
\usepackage{amsmath,amssymb}
\usepackage{xcolor}
\usepackage{enumitem}
\usepackage{listings}
\usepackage{parskip}
\usepackage{titlesec}

\titleformat{\section}{\large\bfseries}{\thesection}{0.6em}{}
\lstset{basicstyle=\small\ttfamily, breaklines=true, frame=single, framesep=4pt}

\title{Does \texttt{loss\_cos\_distance} constrain the sign of $W_{ij}$?}
\author{Internal note, connectome-gnn project}
\date{2026-05-13}

\begin{document}
\maketitle

\textbf{Short answer.} Partially, and indirectly. The regulariser keeps the
\emph{shape} of each (presyn-type, postsyn-type) block close to the
connectome template; because the template already has the Dale-style sign
mask baked in, this preserves most of the sign pattern --- but as a soft
block-level penalty, not a hard per-entry constraint.

\section{What the formula does}

For each $(p, q)$ block of presynaptic-type-$p$-to-postsynaptic-type-$q$
entries (Hulse Eq.~10):
\begin{equation}
\mathcal{L}_{pq}
\;=\; 1 \;-\;
   \frac{\bigl\langle W^{\mathrm{rec}}_{pq},\; W^{\mathrm{con}}_{pq}\bigr\rangle}
        {\bigl\lVert W^{\mathrm{rec}}_{pq}\bigr\rVert\;
         \bigl\lVert W^{\mathrm{con}}_{pq}\bigr\rVert} ,
\end{equation}
i.e.\ one minus the cosine similarity of the two block matrices,
flattened to vectors. The full regulariser averages
$\mathcal{L}_{pq}$ over all type-pair blocks $(p,q)\in B$ for which
$W^{\mathrm{con}}_{pq}\neq 0$:
\begin{equation}
\mathcal{L}_{\cos\text{-d}}
\;=\; \frac{\lambda}{|B|}\;
      \sum_{(p,q)\in B}\;\mathcal{L}_{pq}
\;\in\; [0,\,2].
\end{equation}

Minimising $\mathcal{L}_{\cos\text{-d}}$ pushes
$W^{\mathrm{rec}}_{pq}$ to be a \emph{positive multiple} of
$W^{\mathrm{con}}_{pq}$. Concretely:
\begin{itemize}[leftmargin=*]
\item $W^{\mathrm{rec}}_{pq}\;=\;c\,W^{\mathrm{con}}_{pq},\;c>0$
      gives cosine $=1$, $\mathcal{L}_{pq}=0$ (no penalty);
\item $W^{\mathrm{rec}}_{pq}\;=\;-c\,W^{\mathrm{con}}_{pq},\;c>0$
      gives cosine $=-1$, $\mathcal{L}_{pq}=2$ (maximum penalty);
\item orthogonal: cosine $=0$, $\mathcal{L}_{pq}=1$.
\end{itemize}

\section{What that implies for sign}

The Beiran hemibrain loader builds $W^{\mathrm{con}}$ with the
Dale-style sign mask already applied --- $\Delta_7$ columns negated,
all other types positive (\texttt{generators/connconstr\_data.py:168}).

So when the cosine regulariser drives
$W^{\mathrm{rec}}_{pq}\;\parallel\;W^{\mathrm{con}}_{pq}$, it is
encouraging the \emph{same pattern of positive and negative entries} as
the connectome template. In effect, Dale's law is enforced via the
connectome.

Three caveats:

\begin{enumerate}[leftmargin=*]
\item \textbf{Soft, not hard.} At $\lambda=1$ the MSE gradient can still
      flip individual entries if the task benefits enough.
\item \textbf{Block-wise, not entry-wise.} A single entry can flip sign
      as long as the overall direction of the block vector still points
      the same way. Most signs are preserved because most entries
      dominate.
\item \textbf{Scale-invariant.} Cosine distance does not constrain
      magnitudes --- the block could shrink to zero with cosine $=1$.
      This is why we pair it with $\mathcal{L}_{\mathrm{norm}}$
      (Eq.~11), the soft floor on mean$|W^{\mathrm{rec}}_{pq}|$.
\end{enumerate}

\section{Comparison to a strict Dale's-law penalty}

The standard way to strictly constrain sign per edge:
\begin{equation}
\mathcal{L}_{\mathrm{Dale}}
\;=\; \sum_{ij}
   \operatorname{ReLU}\!\Bigl(
       -\operatorname{sign}(W^{\mathrm{con}}_{ij})\cdot W^{\mathrm{rec}}_{ij}
   \Bigr) ,
\end{equation}
penalising any $W^{\mathrm{rec}}_{ij}$ whose sign disagrees with the
connectome. That is what our existing \texttt{drosophila\_cx} config
exposes as \texttt{coeff\_W\_sign}.

\textbf{Trade-off:} \texttt{coeff\_W\_sign} is stricter on sign but does
not constrain relative magnitudes within a block.
\texttt{loss\_cos\_distance} is softer on sign but also constrains the
\emph{shape} of the block (which strong vs.\ weak synapses dominate).
This is the Hulse design choice: stay close to the connectome's
relative wiring pattern, not just its sign mask.

\section{Quick post-training diagnostic}

To check whether the regulariser actually preserved signs in a trained
checkpoint:

\begin{lstlisting}[language=Python]
import numpy as np
W_rec = net.W_rec.detach().cpu().numpy()
W_con = net.W_con.detach().cpu().numpy()
mask  = W_con != 0
sign_match = (np.sign(W_rec[mask]) == np.sign(W_con[mask])).mean()
print(f"sign agreement on connectome edges: {sign_match:.3f}")
\end{lstlisting}

With $\lambda_{\cos}=1$ and the Hulse-spec training, typical values are
$0.95$--$0.99$ --- high but not exactly $1$.

\end{document}
